\chapter{Threats To Validity}\label{sec:threats}

This chapter reports threats to the validity of the experiments in
Chapter~\ref{s:evaluation}, following the classification proposed by Wohlin
et al.~\cite{wohlin12}.

\section{Internal Validity}

The training and evaluation datasets were both authored and curated by the
same researcher who also designed the pipeline and interpreted the results,
which introduces a risk of confirmation bias in which intents that are known
to work well with the retrieval and prompt design were more likely to be
included or phrased favourably. The regression observed for the first SFT
model (Section~\ref{s:evaluation}) is itself direct evidence of an internal
validity threat that was realised in practice: a prompt-format mismatch
between training and evaluation silently degraded a model's measured
performance and would have been mistaken for a genuine capability regression
had the cause not been specifically investigated. This raises the possibility
that other, undetected train/inference inconsistencies remain in the reported
results. Additionally, the July and August baseline evaluation rounds differ
in more than one variable simultaneously (system prompt, RAG aliases, and
output post-processing all changed together), so the reported 37.9\%
$\rightarrow$ 92.7\% improvement cannot be decomposed into the individual
contribution of each fix.

\section{External Validity}

All evaluated intents target one of seven hand-built Brane packages designed
specifically for this project, covering a curated set of scientific scenarios
(healthcare, genomics, epidemiology, statistics, text analysis, data masking,
datetime). It is not established that the translation accuracy, RAG design,
or cache behaviour reported here generalise to (i)~packages built by other
teams with different documentation conventions, (ii)~substantially larger
package ecosystems, or (iii)~intents drawn from real scientists' requests
rather than researcher-authored examples. The dataset is also comparatively
small (584 unique intents) relative to general-purpose code-generation
benchmarks, which limits confidence that the findings extrapolate to
significantly more diverse workloads.

\section{Construct Validity}

The three evaluation metrics (compile rate, execution rate, output match rate)
are proxies for correctness rather than correctness itself. Execution rate in
particular only certifies that a workflow ran without a runtime error; a
workflow can execute successfully while still computing the scientifically
wrong result if the error is not one that Brane surfaces as a failure (e.g., a
plausible but incorrect statistical computation). Output match rate is a
stricter proxy but compares exact \texttt{stdout} strings, which penalises
workflows that are semantically correct but format their output differently
from the reference (e.g., different rounding or key ordering), and is only
computed over the subset of examples with a stored reference output, which may
not be representative of the full evaluation set. Similarly, the cache
benchmark's ``correct rate'' requires an exact script match, while the
separately reported ``fuzzy correct rate'' (bigram similarity $\geq 0.95$)
is a coarser proxy that may both overstate correctness (superficially similar
but semantically different scripts) and understate it (a correct but
differently formatted script falling below the similarity threshold).

\section{Conclusion Validity}

The 87-example validation split and the single evaluation run per model
configuration limit the statistical strength of the comparisons in
Chapter~\ref{s:evaluation}: no repeated-sampling or significance testing was
performed to quantify the variance introduced by LLM sampling stochasticity,
so reported percentage differences between models (e.g., 92.7\% vs.\ 63.4\%
execution rate) should be read as point estimates from a single trial rather
than as statistically validated effects. The cache threshold sweep is
comparatively more robust, being computed over 1974 paraphrase queries, but
still draws on paraphrases generated by a single LLM (Qwen3.5-4B), which may
itself introduce systematic paraphrasing patterns that do not represent the
full diversity of how a human scientist would reword an intent.